\section{Data}\label{sec:data}

\subsection{Source}

We use Peru's 2024 National Household Survey (Encuesta Nacional de Hogares, ENAHO), a nationally representative repeated cross-sectional survey conducted by the National Institute of Statistics and Informatics (INEI). ENAHO employs a stratified, multistage probability sampling design with primary sampling units (PSUs, identified by the variable \texttt{CONGLOME}) and strata defined by geographic domain and urban/rural classification. We use the survey expansion factor \texttt{FACTOR07} from the Sumaria module throughout all descriptive statistics and regression models to produce population-representative estimates.

We note that the 2024 ENAHO is collected during the post-pandemic recovery period. Peru experienced among the highest COVID-19 excess mortality rates globally, with substantial disruption to health care access and health-seeking behavior. Health expenditure patterns in 2024 may reflect post-pandemic adjustment rather than steady-state behavior, a caveat that limits the temporal generalizability of our findings.

Our analysis draws on two ENAHO modules. Module 400 (Health) provides individual-level information on health care utilization, insurance status, and OOP health expenditure for 110,451 individuals. The Sumaria module provides household-level total consumption (\texttt{GASHOG2D}), which serves as the denominator for our OOP share measure. We merge these modules at the household level using the triple key \texttt{CONGLOME}$\times$\texttt{VIVIENDA}$\times$\texttt{HOGAR}.

\subsection{Variable Construction}

\paragraph{Outcome variable.} Our primary outcome is the OOP health expenditure share, defined as:
\begin{equation}\label{eq:oop_share}
    \text{OOP\_share}_h = \frac{\text{Health\_expenditure}_h}{\max(\text{Total\_consumption}_h, \, \text{Floor})}
\end{equation}
where Health\_expenditure is the sum of consultation fees (\texttt{P41601}), medicine costs (\texttt{P41602}), and other health costs (\texttt{P41603}), aggregated from individual to household level and annualized by multiplying by 13 (52 weeks $\div$ 4-week recall period). Total\_consumption is \texttt{GASHOG2D} from Sumaria, which is already annualized. The floor is set at the 5th percentile of \texttt{GASHOG2D} among households with positive consumption to prevent extreme shares from near-zero denominators. We cap OOP\_share at $[0, 1]$ after applying the floor.

The annualization factor of 13 assumes that 4-week recall expenditure is representative of annual spending. This assumption amplifies measurement error for households experiencing acute episodes during the recall period: a single catastrophic visit is multiplied 13-fold, potentially overstating annual OOP. We report sensitivity to an alternative $\times 12$ annualization in Section~\ref{sec:robustness} and note that this amplification is the dominant source of measurement error in the outcome variable. As an additional consideration, the \texttt{P41603} category (``other health costs'') captures heterogeneous items including hospitalization, transport, dental, and optical expenses, whose high within-category variance may disproportionately drive upper-quantile estimates.

\paragraph{Catastrophic health expenditure.} We construct binary CHE indicators at three thresholds: $\text{CHE}_{10} = \mathbf{1}[\text{OOP\_share} > 0.10]$, $\text{CHE}_{25} = \mathbf{1}[\text{OOP\_share} > 0.25]$, and $\text{CHE}_{40} = \mathbf{1}[\text{OOP\_share} > 0.40]$.

\paragraph{Insurance status.} We construct mutually exclusive insurance categories from Module 400: (i) SIS (\texttt{P4191}$=1$ and \texttt{P4192}$\neq 1$), (ii) EsSalud (\texttt{P4192}$=1$, with or without SIS; EsSalud takes precedence in overlaps), and (iii) Uninsured (neither SIS nor EsSalud). The household-level insurance variable is assigned based on the household head's coverage.

We acknowledge an important data limitation: our EsSalud sample comprises only $N = 648$ household heads (1.9\% of the sample), which is substantially below the expected 30--35\% national EsSalud coverage rate. This discrepancy likely reflects a variable construction issue in the P4192 coding within Module 400---possibly related to how multiple insurance affiliations are recorded at the individual level, or to the fact that Module 400 captures a different population than the insurance-specific modules used in official coverage statistics. We were unable to resolve this discrepancy within the current analysis. Consequently, \textbf{EsSalud results throughout this paper should be interpreted with extreme caution}. The small cell sizes ($N \approx 8$--130 per quintile) render EsSalud--quintile interaction estimates imprecise and potentially unreliable. We retain EsSalud in the analysis for completeness but emphasize that the primary substantive conclusions rest on the SIS--quintile comparison.

An additional limitation of the head-based insurance assignment is that it excludes cases where the head is uninsured but dependents have SIS coverage, or vice versa. A household-level measure based on any-member coverage may better characterize financial protection. We report sensitivity to alternative aggregation rules in Section~\ref{sec:robustness}.

\paragraph{Consumption quintiles.} Quintiles are constructed from survey-weighted \texttt{GASHOG2D} using \texttt{FACTOR07}-weighted quantile breaks. We note a mechanical circularity: \texttt{GASHOG2D} is a comprehensive consumption aggregate that, by construction, includes health expenditures (Module 400 flows into Sumaria). Households with high OOP spending are thus assigned to higher quintiles partly because of their health spending, inducing a positive correlation between OOP share and quintile rank that inflates the quintile gradient. The standard correction---constructing quintiles from consumption net of health expenditure---would require subtracting annualized health spending from \texttt{GASHOG2D} before computing quintile breaks. This circularity is a fundamental limitation of the current implementation: the quintile--OOP gradient reported in our results is partly mechanical rather than behavioral. Net-of-health quintiles---constructed from $\text{GASHOG2D} - \text{annualized health expenditure}$---are required as the primary specification to isolate the behavioral gradient. We were unable to implement this correction in the current version and flag this as a priority for reanalysis. All quintile-related results, including the insurance--quintile interactions that form the core of our analysis, should be interpreted as upper bounds on the true behavioral association.

\paragraph{Controls.} Household head characteristics include age (\texttt{P208A}), sex (\texttt{P207}), and years of education (\texttt{P301A}). Household-level variables include household size, number of children under 5, number of elderly members (65+), rural residence (\texttt{AREA}), and presence of any chronic illness among household members. We include region fixed effects based on INEI's 24 administrative departments (derived from the first two digits of \texttt{UBIGEO}).

\subsection{Sample and Descriptive Statistics}

After merging modules and aggregating to household level, our analysis sample comprises 33,691 households. Table~\ref{tab:sumstats} reports survey-weighted summary statistics. The mean OOP share is 2.7\% ($\text{SD} = 9.2\%$), with a median of 0.3\%, reflecting the pronounced right skew of health spending. A substantial 41.1\% of households report zero OOP health expenditure---a feature with critical implications for quantile regression, as discussed in Section~\ref{sec:strategy}. CHE prevalence is 5.1\% at the 10\% threshold, 1.3\% at 25\%, and 0.8\% at 40\%.

\begin{table}[htbp]
\centering
\caption{Summary Statistics (Survey-Weighted)}
\label{tab:sumstats}
\small
\begin{tabular}{lrrrrrrr}
\toprule
Variable & N & Mean & SD & Min & Median & Max \\
\midrule
  oop\_share & 33,691 & 0.027 & 0.092 & 0.000 & 0.003 & 1.000 \\
  health\_expend & 33,691 & 11415.231 & 136853.281 & 0.000 & 65.000 & 3910306.400 \\
  total\_consumption & 33,691 & 31285.282 & 22758.677 & 825.153 & 23945.350 & 308555.781 \\
  zero\_oop & 33,691 & 0.411 & 0.492 & 0.000 & 0.000 & 1.000 \\
  che\_10pct & 33,691 & 0.051 & 0.220 & 0.000 & 0.000 & 1.000 \\
  che\_25pct & 33,691 & 0.013 & 0.113 & 0.000 & 0.000 & 1.000 \\
  che\_40pct & 33,691 & 0.008 & 0.089 & 0.000 & 0.000 & 1.000 \\
  age\_head & 33,691 & 53.430 & 15.586 & 15.000 & 53.000 & 98.000 \\
  female\_head & 33,691 & 0.373 & 0.484 & 0.000 & 0.000 & 1.000 \\
  educ\_years & 33,691 & 9.885 & 5.514 & 0.000 & 11.000 & 20.000 \\
  hh\_size & 33,691 & 3.344 & 1.760 & 1.000 & 3.000 & 17.000 \\
  children\_u5 & 33,691 & 0.193 & 0.448 & 0.000 & 0.000 & 5.000 \\
  elderly\_65 & 33,691 & 0.409 & 0.672 & 0.000 & 0.000 & 5.000 \\
  rural & 33,691 & 0.205 & 0.403 & 0.000 & 0.000 & 1.000 \\
  chronic\_any & 33,691 & 0.163 & 0.369 & 0.000 & 0.000 & 1.000 \\
  hospitalized & 33,691 & 0.000 & 0.000 & 0.000 & 0.000 & 0.000 \\
\bottomrule
\end{tabular}
\begin{minipage}{\textwidth}\footnotesize
\textit{Notes:} ENAHO 2024. Survey-weighted (FACTOR07). Unit: household.
\end{minipage}
\end{table}

Table~\ref{tab:balance} presents descriptive statistics by insurance group. SIS enrollees ($N = 11{,}189$, 33.2\% of the sample) have a mean OOP share of 3.0\% and a zero-OOP rate of 37.3\%. EsSalud members ($N = 648$, 1.9\%) show a slightly higher mean OOP share (3.4\%) and zero-OOP rate (40.9\%). Uninsured households ($N = 1{,}165$, 3.5\%) have the lowest mean OOP share (2.5\%) but the highest zero-OOP rate (48.7\%), consistent with care avoidance among the uninsured. SIS households are more rural (5.9\% vs.\ 1.0\% for EsSalud), have more children under 5, and lower education levels, reflecting the means-tested targeting of SIS toward poorer households.

\begin{table}[htbp]
\centering
\caption{Descriptive Statistics by Insurance Group (Survey-Weighted)}
\label{tab:balance}
\small
\begin{tabular}{lcccccc}
\toprule
& \multicolumn{2}{c}{SIS} & \multicolumn{2}{c}{EsSalud} & \multicolumn{2}{c}{Uninsured} \\
\cmidrule(lr){2-3} \cmidrule(lr){4-5} \cmidrule(lr){6-7}
Variable & Mean & SD & Mean & SD & Mean & SD \\
\midrule
  oop\_share & 0.030 & 0.104 & 0.034 & 0.136 & 0.025 & 0.074 \\
  zero\_oop & 0.373 & 0.484 & 0.409 & 0.492 & 0.487 & 0.500 \\
  health\_expend & 16440.741 & 172406.483 & 30701.373 & 235556.493 & 3599.961 & 61523.844 \\
  age\_head & 55.569 & 15.172 & 54.861 & 14.727 & 52.019 & 15.177 \\
  female\_head & 0.380 & 0.485 & 0.379 & 0.485 & 0.329 & 0.470 \\
  educ\_years & 11.789 & 5.638 & 13.068 & 7.120 & 9.867 & 5.230 \\
  hh\_size & 3.555 & 1.742 & 3.601 & 1.637 & 2.039 & 1.317 \\
  children\_u5 & 0.171 & 0.425 & 0.136 & 0.368 & 0.036 & 0.202 \\
  elderly\_65 & 0.506 & 0.738 & 0.494 & 0.732 & 0.260 & 0.515 \\
  rural & 0.059 & 0.236 & 0.010 & 0.099 & 0.161 & 0.367 \\
  chronic\_any & 0.172 & 0.378 & 0.153 & 0.360 & 0.084 & 0.277 \\
  hospitalized & 0.000 & 0.000 & 0.000 & 0.000 & 0.000 & 0.000 \\
\midrule
  N & \multicolumn{2}{c}{11,189} & \multicolumn{2}{c}{648} & \multicolumn{2}{c}{1,165} \\
\bottomrule
\end{tabular}
\begin{minipage}{\textwidth}\footnotesize
\textit{Notes:} ENAHO 2024. Survey-weighted means (FACTOR07).
\end{minipage}
\end{table}

\subsection{Zero-OOP Mass and Implications for Quantile Regression}

The 44.1\% zero-OOP rate has direct methodological implications. At quantiles below the zero-mass fraction, the conditional quantile of OOP share is identically zero for most covariate configurations, rendering standard Koenker--Bassett quantile regression coefficients uninformative \citep{MachadoSilva2005}. At $\tau = 0.10$ and $\tau = 0.25$, estimated coefficients reflect movements between zero and near-zero values rather than meaningful variation in spending intensity. At $\tau = 0.50$, the conditional quantile is positive for most but not all covariate cells---particularly low-income, low-utilization subgroups may still have degenerate medians.

We address this zero-mass problem through two complementary strategies. First, we restrict substantive interpretation of CQR results to $\tau \geq 0.50$, where conditional quantiles are positive for most household types, and treat $\tau = 0.10$ and $\tau = 0.25$ estimates as uninformative lower-bound specifications reported for completeness. Second, we complement the CQR analysis with a two-part model that separately identifies the extensive margin (probability of any OOP spending via probit) and the intensive margin (QR on the positive-OOP subsample), following \citet{Manning1987}. We note that censored quantile regression \citep{Powell1986, Chernozhukov2002} offers an alternative approach that treats zero as a censoring point and produces consistent estimates at all quantiles; this remains a direction for future work.
